% Nota de k_por_modelo. GERADA por flim_al/tabelas.py.
\footnotesize Uma linha por (critério, K); cada coluna de decoder é o F$\beta$ naquele decoder. \texttt{flim\_paper} é o braço do artigo --- as imagens fixas que os autores escolheram, sem seleção automática. \textbf{Todos os braços usam o mesmo gerador de marker sintético}, inclusive o do artigo: marker real só existe para 31 imagens, e misturar as duas origens faria o $\Delta$ medir quem desenhou o traço em vez da seleção. A comparação é internamente válida e \textbf{não} se compara com números produzidos a partir de markers reais. \texttt{treino} é estimar os kernels por k-means --- o que o FLIM promete ser barato; \texttt{teste} é rodar os sete decoders no conjunto de avaliação. Ele cresce com K porque o encoder cresce: a arquitetura pede 200 kernels por camada, mas o k-means só produz tantos quantos os patches dos markers permitem --- medido, 54/51/48/48 com uma imagem contra 200/200/200/200 com oito. Com poucas anotações a rede não é só menos treinada, é menor. Tempo com \texttt{~} é estimado a partir da soma treino+avaliação, que era tudo que o registro guardava antes; sem \texttt{~} é medido. $\Delta$ é contra o artigo no MESMO K, e fica vazio em K acima de cinco porque o braço do artigo tem cinco imagens e não existe referência para comparar. Todas as linhas vêm de uma única campanha (\texttt{tabela\_k\_por\_modelo/semente0/pool40/teste60}): campanhas diferentes usam conjuntos de teste diferentes e não se misturam na mesma linha.
\normalsize
